%% notes-tex/031-inhibitor-tolerance-data/sections/1-question.tex

\section{The question\secstatus{todo}}
\label{sec:question}

A strain-design model for industrial hosts has to predict how a genotype responds to
the inhibitors a real feedstock carries. Lignocellulosic hydrolysate brings furans
(furfural, 5-hydroxymethylfurfural), weak acids (acetic, levulinic, formic), phenolics
from lignin (vanillin, syringaldehyde, ferulic and p-coumaric acid), and the ionic
liquids left by pretreatment; a producing strain also has to tolerate its own product,
which for the isobutanol work is isobutanol. Vanacloig-Pedros et al.\ 2022
(\texttt{10.1093/femsyr/foac036}) screened a barcoded deletion library against exactly
that panel, and it is the only genome-wide dataset in the served graph that does. The
question for a model is what other data can help it, and the candidate is the only
other genome-wide environment-response compendium of comparable size, Hillenmeyer et
al.\ 2008 (\texttt{10.1126/science.1150021}), in its homozygous (HOM) and heterozygous
(HET) arms.

Every script named below lives in
\file{experiments/031-env-chemgen-inhibitor-tolerance/scripts/}. Three things have to
be known before a model is trained on the pair. First, the axes
on which the datasets differ, because an environment encoding has to carry every axis
that changes between them and a cell encoding has to carry every axis that changes
inside the strain. Second, how well each dataset agrees with itself, because replicate
agreement is the ceiling on any prediction and it is not the same for every compound.
Third, how much response structure the two datasets share, because a partner that
measures different biology adds parameters and no signal.

\notesec{Vocabulary}
A \term{condition} is one dosed compound, or one physical factor or temperature when no
compound is dosed; Hillenmeyer's several doses and generation counts of one compound are
one condition here. The \term{response} is the served number for a strain in a
condition, oriented throughout so that negative is a fitness defect. A \term{hit} is a
gene in the bottom 5\% of responses in a condition. The \term{reliability index} of a
condition is $1 - \overline{\mathrm{SE}^2} / \mathrm{Var}(\text{response})$ over its
genes, the share of the across-gene variance that the served replicate noise cannot
produce. \term{Replicate agreement} is the Spearman correlation over genes between two
raw replicates of a condition, and \term{hit agreement} the Jaccard index of their hit
sets. The \term{cross Spearman} of a Vanacloig condition and a Hillenmeyer condition is
the Spearman correlation of their responses over the genes both screened. An
\term{exact match} is a compound both datasets dose under the same InChIKey;
\term{chemical similarity} is Tanimoto similarity of fingerprints or cosine similarity of
a learned embedding between two compounds that are not the same molecule.

\begin{keybox}
\textbf{Measured.} Every axis of the three datasets from their served records
(Section~\ref{sec:axes}); replication, served error and per-condition ceilings
(Section~\ref{sec:noise}); the shared response structure and the transfer of every
exact-match compound (Section~\ref{sec:structure}); the coverage of every molecular
encoder over the three compound sets and the relation between chemical similarity and
response similarity (Sections~\ref{sec:encoders} and \ref{sec:chemistry}).

\textbf{Not measured.} No model has been trained. The linear and nearest-neighbor
baselines that turn these ceilings into scores are Section~\ref{sec:next}.
\end{keybox}
